\documentclass[10pt,a4paper,twoside]{article}
\usepackage[dutch]{babel}
%laad de pakketten nodig om wiskunde weer te geven :
\usepackage{amsmath,amssymb,amsfonts,textcomp}
%laad de pakketten voor figuren :
\usepackage{graphicx}
\usepackage{float,flafter}
\usepackage{hyperref}
\usepackage{inputenc}
%zet de bladspiegel :
\setlength\paperwidth{20.999cm}\setlength\paperheight{29.699cm}\setlength\voffset{-1in}\setlength\hoffset{-1in}\setlength\topmargin{1.499cm}\setlength\headheight{12pt}\setlength\headsep{0cm}\setlength\footskip{1.131cm}\setlength\textheight{25cm}\setlength\oddsidemargin{2.499cm}\setlength\textwidth{15.999cm}

\begin{document}
\begin{center}
\hrule

\vspace{.4cm}
{\bf {\Huge Mijn Verslag}}
\vspace{.2cm}
\end{center}
{\bf Mijn Naam}  \\
{\bf Eerste Bachelor Fysica en Sterrenkunde} \hspace{\fill} Groep {\bf Z}\\
Opstelling {\bf 0 } \hspace{\fill}  1 januari 2025 \\
\hrule
%\bf genereert vette letters, \large \Large \huge \Huge \tiny \small ... verschillende lettergroottes, met \em, \it, \sl krijg je cursieve tekst
%Je moet accolades gebruiken om aan te geven waarop je het commando precies wil laten inwerken.
% commentaar zet je in de tex-file met een %-tekentje voor.


\section{Inleiding}

Hier komt de inleidende tekst. Zorg dat deze duidelijk is, niet te lang en schrijf geen inleidingsnota's over~!
%de ~ zorgt ervoor dat het ! met een spatie aan de tekst wordt geplakt en niet naar de volgende lijn verhuist.  Omdat latex na elk . automatisch een dubbele spatie invoegt, kan je ~ ook gebruiken om te voorkomen dat je na een afkorting bvb.~een dubbele spatie krijgt.

%1 open lijn wordt door latex gewoon als spatie gezien, met twee linefeeds  begint een nieuwe paragraaf.  Als je wil voorkomen dat de tekst inspringt, kan je het commando 
oindent gebruiken.  \\ springt naar een volgende lijn, 
ewpage naar een volgend blad.


%\begin{figure}
%\centering
%\includegraphics[width=0.7\linewidth]{figje}
%\caption{}
%\label{fig:figje}
%\end{figure}
\section{Experimentele methode}
In deze paragraaf worden de experimentele opstelling en methode besproken.


\section{Meetresultaten}
%Met het commando \sectie{naamtitel} maak je een nieuwe titel aan.
De meetresultaten worden samengevat.  

\subsection{Meting van 1... en bespreking}
Eventueel kan je een paragraaf op deze manier onderverdelen in verschillende subsecties.

%hier werd verwezen naar de tabel die als \label (naam} 'tabel1' kreeg.




\subsection{Meting 2}
... enzoverder.

\section{Bespreking}
Hier komt de bespreking als die niet bij de meetresultaten werd gezet. 
Opsommingen kan je invoeren als~:
\begin{enumerate}
\item 
eerste puntje
\item
tweede puntje
\item
en zelfs genest zoals bijvoorbeeld
\begin{itemize}
\item eerste subpuntje
\item ....
\end{itemize}
\end{enumerate}
%met enumerate genereer je een opsomming, itemize maakt verschillende puntjes

\section{Besluit}
In deze paragraaf worden de besluiten geformuleerd.

\end{document}